% ---
% title: 操作系统期末真题（一）解析
% date: 2026-09-22
% author: Crystal-Sky
% tags: 操作系统, 期末复习, 真题解析
% summary: 操作系统期末真题（一）题目与逐题解析：磁盘空闲块管理、死锁判断、虚拟页式地址转换、并发共享变量、文件系统目录查找与调度、同步算法设计。
% slug: os-final-exam-1
% course: 操作系统
% course_slug: operating-systems
% chapter: 期末真题
% lesson_order: 1
% lesson_type: 真题解析
% challenge: false
% inline_answers: true
% ---
\documentclass[UTF8,12pt]{ctexart}
\usepackage[a4paper,margin=2.2cm]{geometry}
\usepackage{amsmath,amssymb}
\usepackage{array,booktabs,tabularx,longtable}
\usepackage{enumitem}
\usepackage{listings}
\usepackage{xcolor}
\usepackage{graphicx}
\usepackage[colorlinks=true,linkcolor=blue,urlcolor=blue]{hyperref}
\setlength{\parindent}{2em}
\setlength{\parskip}{0.35em}
\linespread{1.25}
\setlist{nosep,leftmargin=2em}

\newenvironment{question}{\par\medskip\noindent\textbf{【题目】}\par\smallskip}{\par\medskip}
\newenvironment{explanation}{\par\noindent\textbf{【解析】}\par\smallskip}{\par\medskip}
\newenvironment{answer}{\par\noindent\textbf{【答案】}\par\smallskip}{\par}

\lstset{
  basicstyle=\ttfamily\small,
  columns=fullflexible,
  keepspaces=true,
  showstringspaces=false,
  frame=single,
  breaklines=true
}

\begin{document}
\section*{一、（24分）简要回答下列问题。}

\subsection*{1.（6分）磁盘空闲存储空间管理}
\begin{question}
1.（6分）对于磁盘空闲存储空间管理，下列 2 种方法哪一种好？为什么？

（1）每个空闲块中指出下一个空闲块的块号。

（2）采用多个空闲块构成的链表存放空闲块号。
\end{question}

\begin{answer}
通常第（2）种方法更好。原因是它一次可以取得多个空闲块号，减少访问磁盘的次数，从而降低空闲块分配、回收时的 I/O 开销。
\end{answer}

\begin{explanation}
第（1）种方法本质上是把所有空闲块逐块串成链表：每访问一个链表结点，通常都要访问一个磁盘块，若需要连续取得较多空闲块，磁盘 I/O 次数较多。

第（2）种方法把多个空闲块号集中存放在一个空闲块中，再通过链接把这些“成组”的空闲块号组织起来。这样一次读入一个管理块就能获得多个空闲块号，可以明显减少为寻找空闲块而进行的磁盘访问次数，分配和回收的效率更高。这就是常见的成组链接思想。
\end{explanation}

\subsection*{2.（6分）死锁判断}
\begin{question}
2.（6分）下面关于死锁问题的叙述，是否正确？为什么？

（1）若 A、B、C、D 四个进程都处于阻塞状态，则 A、B、C、D 构成的进程集合出现死锁。

（2）若 A、B、C、D 四个进程构成的集合出现死锁，则 A、B、C、D 中的每个进程都一定占有资源。

（3）若系统处于不安全状态，则一定会导致死锁。
\end{question}

\begin{answer}
（1）错误；（2）正确；（3）错误。
\end{answer}

\begin{explanation}
\textbf{（1）错误。}“阻塞”不等于“死锁”。进程可能只是等待 I/O 完成、等待某个事件发生，等待条件满足后仍能继续执行。因此四个进程同时阻塞，并不能推出它们之间存在循环等待，也不能推出发生死锁。

\textbf{（2）正确（按资源死锁的通常模型理解）。}构成死锁集合的进程之间形成循环等待：每个进程至少占有某些资源，同时等待集合中其它进程占有的资源。如果其中某个进程完全不占有任何资源，它就不能成为该循环中阻塞其它进程的环节。

\textbf{（3）错误。}不安全状态表示系统已经不能保证所有进程都存在一个安全的完成次序，但这只是“有可能发生死锁”，并不表示当前已经死锁，也不表示以后一定会死锁。如果后续进程对资源的实际需求较小，仍可能全部顺利结束。
\end{explanation}

\subsection*{3.（6分）共享设备}
\begin{question}
3.（6分）I/O 设备按共享属性分类可分为独占设备和共享设备。举一个共享设备的例子，针对你所举的例子，说明共享的含义，给出一种设备分配方法，并说明你的方法存在的不足。
\end{question}

\begin{answer}
例：磁盘。多个进程可以共同提交磁盘访问请求，由操作系统统一排队调度。可采用 FCFS 方法；不足是没有优化磁头移动，平均寻道时间可能较长，性能较低。
\end{answer}

\begin{explanation}
以\textbf{磁盘}为例。

磁盘属于共享设备。这里的“共享”不是指多个进程在同一瞬间同时驱动磁头，而是指系统不把整个磁盘在一段较长时间内独占地分配给某一个进程。多个进程都可以提交磁盘 I/O 请求，操作系统把这些请求排入请求队列，再按某种磁盘调度算法逐个处理。

一种简单的分配方法是\textbf{先来先服务（FCFS）}：磁盘请求按照到达先后进入队列，设备空闲时总是选择队首请求执行。

这种方法的优点是简单、比较公平，也不容易出现请求长期得不到服务的问题；不足是没有考虑磁头当前位置，可能造成磁头来回大距离移动，使平均寻道时间较大，磁盘吞吐量较低。
\end{explanation}

\subsection*{4.（6分）内存管理设计目标的冲突}
\begin{question}
4.（6分）在设计操作系统的内存管理功能时，需要考虑多个设计目标，例如内存利用率高、便于实现内存共享等，而有些设计目标之间往往存在冲突，需要权衡。举 2 个与内存管理相关的设计目标（不限于本题所列的设计目标）存在冲突的例子，分别解释可能冲突的情况，并简要阐述在设计时采取什么策略对相互冲突的设计目标进行权衡。
\end{question}

\begin{answer}
示例答案：

1. 大页面可减小页表和管理开销，但会增加内部碎片；可选取适中的页大小或采用多种页大小。

2. 给每个进程较少页框可提高多道程序度，但可能增加缺页甚至产生抖动；可通过工作集或缺页率控制动态调整页框分配。
\end{answer}

\begin{explanation}
可以举出很多组合理的冲突目标，下面给出两组典型例子。

\textbf{例 1：页大小带来的“管理开销”与“内存利用率”之间的冲突。}

页越大，一个进程需要的页表项越少，页表占用空间小，地址转换和换页时的管理开销也往往较低；但是页越大，最后一页没有用满所造成的内部碎片通常越严重，内存利用率会下降。页越小则相反。

设计时通常选择一个折中的页大小，或者支持多种页大小：普通内存使用较小页面，需要大容量、连续映射时使用大页。

\textbf{例 2：提高多道程序度与降低缺页率之间的冲突。}

如果给每个进程分配较少页框，就可以让更多进程同时驻留内存，提高多道程序度；但每个进程可用的页框太少时，会频繁缺页，甚至发生抖动。反之，给每个进程较多页框可以降低缺页率，但同时驻留内存的进程数量会减少。

设计时可以采用工作集模型、缺页率（PFF）控制等方法，根据进程实际内存需求动态调整页框数量，在多道程序度和缺页率之间进行折中。
\end{explanation}

\section*{二、（18分）虚拟页式存储管理。}

\begin{question}
某计算机系统采用虚拟页式存储管理，虚拟地址为 32 位，页的大小为 4096 个字节。采用请求调页、固定分配和局部置换策略，页的置换以及快表项的置换均采用 FIFO（First-In-First-Out）算法，页表在内存中。产生缺页中断时，系统将缺页调入内存，更新页表（不更新快表），返回到产生缺页中断的指令处重新执行。

1. 假定为进程 P 分配了 3 个物理块（或称页框、页帧，pageframe），快表长度为 2，进程 P 依次有页号为 10H、100H 和 60H 的页面调入内存中。此时，进程 P 被调度执行，快表为空，依次访问下列虚拟地址（假定这些虚拟地址均在进程 P 的虚拟地址空间内）：

（1）100860H　（2）20E60H　（3）100920H　（4）60100H　（5）100600H

得到这些虚拟地址对应的物理地址各需访问内存几次？说明每次访问内存的功能。

2. 指出该内存管理机制存在的一种不足。针对这种不足，给出一种改进方法，并解释理由。
\end{question}

\begin{answer}
第 1 问：按“得到物理地址”统计页表主存访问次数，答案为
\[
\boxed{1,\;2,\;0,\;1,\;1}.
\]
若把最终访问目标单元的一次主存访问也计入，则为
\[
\boxed{2,\;3,\;1,\;2,\;2}.
\]

第 2 问：缺页置换后只更新页表、不处理快表，可能留下失效的快表项。改进方法是在页面换出时使对应 TLB 项失效，并在需要时装入新的有效映射，保证页表和快表一致。
\end{answer}

\begin{explanation}
页大小为 $4096=2^{12}=\mathrm{1000H}$ 字节，因此一个虚拟地址的低 12 位（即低 3 个十六进制位）是页内偏移，高位是页号。

所以五个地址对应的页号分别为：
\[
\mathrm{100H},\quad \mathrm{20H},\quad \mathrm{100H},\quad \mathrm{60H},\quad \mathrm{100H}.
\]

初始状态：
\[
\text{内存中的页面（FIFO 先后次序）}:\quad 10H\rightarrow100H\rightarrow60H,
\]
\[
\text{快表}:\quad \varnothing.
\]

下面只统计“\textbf{为了得到物理地址而访问内存}”的次数，也就是访问内存中的页表的次数；访问快表本身不算访问主存。

\begin{center}
\begin{tabularx}{\textwidth}{>{\centering\arraybackslash}p{1.1cm} >{\centering\arraybackslash}p{2.3cm} >{\centering\arraybackslash}p{1.6cm} X >{\centering\arraybackslash}p{1.8cm}}
\toprule
次序 & 虚拟地址 & 页号 & 过程 & 为得到物理地址访问内存次数\\
\midrule
1 & 100860H & 100H & 快表为空，快表未命中。访问内存中的页表，发现 100H 已在内存，取得页框号，并把 100H 的页表项装入快表。 & 1\\
\addlinespace
2 & 20E60H & 20H & 快表未命中。第一次访问页表发现 20H 不在内存，发生缺页。按页面 FIFO 淘汰最早进入内存的 10H，调入 20H 并更新页表，但不更新快表。指令重新执行后快表仍未命中，因此第二次访问页表，取得 20H 的页框号，再把该项装入快表。 & 2\\
\addlinespace
3 & 100920H & 100H & 100H 已在快表中，快表命中，直接得到页框号，不必访问内存中的页表。 & 0\\
\addlinespace
4 & 60100H & 60H & 快表未命中，访问页表得到 60H 的页框号。把 60H 装入长度为 2 的快表；按 FIFO 淘汰最早进入快表的 100H。 & 1\\
\addlinespace
5 & 100600H & 100H & 此时 100H 仍在内存，但已经被快表 FIFO 淘汰，因此快表未命中。访问页表取得其页框号。 & 1\\
\bottomrule
\end{tabularx}
\end{center}

因此，按题目“\textbf{得到物理地址}”的字面含义，五次访问分别需要：
\[
\boxed{1,\;2,\;0,\;1,\;1}
\]
次主存访问。

如果课程讲义把“取得物理地址后真正访问目标数据/指令的那一次主存访问”也一起计算，则每项还要再加 1，得到：
\[
\boxed{2,\;3,\;1,\;2,\;2}.
\]
这种口径统计的是“完成一次虚拟地址访问”的主存访问次数，而不是单纯“得到物理地址”的次数。

\medskip
\textbf{第 2 问：}

题目中特别规定：发生缺页时“更新页表（不更新快表）”。这会带来一个重要问题：如果被淘汰的页面恰好还存在于快表中，那么它原来的快表项会变成\textbf{失效的旧映射}。之后再次访问该页时，若仍然命中这个旧快表项，就可能得到已经被其它页面占用的页框号，造成错误的地址转换。

改进方法是在页面被换出、页表项发生变化时，同时\textbf{使对应快表项失效（invalidate）}；也可以在缺页处理结束时将新页的有效映射装入快表。这样可以保证快表与页表的一致性，避免使用陈旧映射。
\end{explanation}

\section*{三、（12分）并发执行与共享变量。}

\begin{question}
设有两个程序 P1 与 P2，其执行的代码分别如下。

程序 P1：
\begin{lstlisting}
int x;
for (i = 0; i < 3; i++)
    x = x + 1;
\end{lstlisting}

程序 P2：
\begin{lstlisting}
int x;
for (i = 0; i < 2; i++)
    x = x - 1;
\end{lstlisting}

其中，x 是共享变量，初值为 0。现启动 4 个进程（2 个进程对应程序 P1，2 个进程对应程序 P2）并发执行，请给出这 4 个进程均执行完成后所有可能的 x 值。
\end{question}

\begin{answer}
\[
\boxed{x\in\{-4,-3,-2,-1,0,1,2,3,4,5,6\}}
\]
\end{answer}

\begin{explanation}
关键点在于：
\[
x=x+1\quad\text{和}\quad x=x-1
\]
并不是不可分割的原子操作。一次“加 1”至少可以抽象为：
\[
\text{读 }x\;\longrightarrow\;\text{在寄存器中加 1}\;\longrightarrow\;\text{写回 }x,
\]
“减 1”同理。因此，一个进程读到旧值以后，在它写回之前，别的进程可能已经修改了 $x$；随后旧的计算结果又可能覆盖新的结果，产生\textbf{丢失更新}。

两个 P1 一共执行 6 次“加 1”，两个 P2 一共执行 4 次“减 1”。

若完全串行，不发生丢失更新，则最终值为
\[
0+6-4=2.
\]
但并发交叉执行时，某些加法或减法结果可以被其它进程基于旧值的写操作覆盖。

最大可能值为 $6$：可以通过合适的交叉，使 4 次减法的影响最终都被后续加法的写回覆盖，而 6 次加法形成最终累计结果。

最小可能值为 $-4$：同理，可以使 6 次加法的影响最终被后续减法覆盖，而 4 次减法形成最终累计结果。

在 $-4$ 到 $6$ 之间的每一个整数，都可以通过安排不同的读、写交叉次序得到。因此所有可能值为连续的整数区间：
\[
\boxed{\{-4,-3,-2,-1,0,1,2,3,4,5,6\}}.
\]

注意：如果把整条 C 语句错误地当成原子操作，则只能得到 $2$；本题考查的正是共享变量读--改--写操作在并发情况下的竞态。
\end{explanation}

\section*{四、（16分）文件系统与多级目录。}

\begin{question}
某文件系统采用多级目录结构，目录当作文件。目录项只包含文件名和 I 节点号，文件地址（磁盘物理块号）等其它文件属性放在 I 节点中。假定磁盘逻辑块与物理块的大小均为 1K 字节，磁盘物理块（以下简称磁盘块）号用 4 个字节表示，根据 I 节点号可计算出其 I 节点所在的磁盘块号，根据磁盘块号可确定磁盘物理地址。目录项中的文件地址由 10 个地址组成，其中 8 个地址是直接地址，1 个地址是一次间接地址，1 个地址项是二次间接地址，每个地址只包含磁盘块号。设文件 f.dat 存在于目录 /zhang/data 下。

1. 设每个目录文件所占空间不超过 2 个磁盘块，每个 I 节点所占空间不超过 1 个磁盘块，根目录文件的 I 节点号已知。为了获得文件 f.dat 的地址，最少需要读多少个磁盘块？最多需要读多少个磁盘块？要求给出计算过程。

2. 你认为如何提高根据文件名（文件路径名）获得文件地址的速度？请给出一种方法，并解释具体理由。

3. 在获得文件 f.dat 的地址后，若要读文件 f.dat 的第 269K 个字节（字节编号从 0 开始），需要读几个磁盘块？要求给出计算过程。
\end{question}

\begin{answer}
1. 最少：\(\boxed{7}\) 个磁盘块；最多：\(\boxed{10}\) 个磁盘块。

2. 可使用目录项/路径名缓存（同时可配合 I 节点缓存），减少逐级目录查找所需的磁盘 I/O。

3. 第 269K 个字节位于逻辑块 269，已经进入二次间接范围，需要读取“一级间接块 + 二级间接块 + 数据块”，共 \(\boxed{3}\) 个磁盘块。
\end{answer}

\begin{explanation}
\textbf{第 1 问：}

路径为
\[
/\;\longrightarrow\;zhang\;\longrightarrow\;data\;\longrightarrow\;f.dat.
\]

已知的是\textbf{根目录文件的 I 节点号}，但要知道根目录文件存放在哪些磁盘块，仍然要先根据 I 节点号读根目录的 I 节点。

沿路径逐级查找时，需要进行：
\begin{enumerate}
  \item 读根目录 I 节点；
  \item 读根目录文件，在其中找到 zhang 的 I 节点号；
  \item 读 zhang 的 I 节点；
  \item 读 zhang 目录文件，在其中找到 data 的 I 节点号；
  \item 读 data 的 I 节点；
  \item 读 data 目录文件，在其中找到 f.dat 的 I 节点号；
  \item 读 f.dat 的 I 节点，从而获得文件地址。
\end{enumerate}

共有 4 个 I 节点需要读取：根目录、zhang、data、f.dat。每个 I 节点最多占 1 个磁盘块，因此 I 节点读取共需 4 个磁盘块。

需要搜索的目录文件有 3 个：根目录、zhang、data。

最理想时，每一级所需目录项都在读到的第 1 个目录块中，因此目录文件共读 3 个块：
\[
4+3=\boxed{7}\text{ 个磁盘块}.
\]

最坏时，每个目录文件都要读满 2 个块才能找到目标目录项，因此目录文件共读 $3\times2=6$ 个块：
\[
4+6=\boxed{10}\text{ 个磁盘块}.
\]

\medskip
\textbf{第 2 问：}

可以建立\textbf{目录项缓存（dentry cache）/路径名缓存}，把最近访问过的“父目录 I 节点号 + 文件名 $\rightarrow$ 子文件 I 节点号”的映射保存在内存中，同时缓存常用 I 节点。

以后再次访问相同或相近路径时，很多路径分量可以直接在内存缓存中命中，不必每一级都重新读目录文件和 I 节点，因此能够显著减少磁盘 I/O，提高路径名解析速度。

\medskip
\textbf{第 3 问：}

一个磁盘块大小为 $1K$，因此第 $269K$ 个字节的逻辑块号为
\[
\frac{269K}{1K}=269.
\]
字节编号从 0 开始，所以直接地址覆盖逻辑块号
\[
0\sim7,
\]
共 8 块。

一个间接地址块大小为 $1K=1024$ 字节，每个磁盘块号占 4 字节，因此一次间接块可以保存
\[
\frac{1024}{4}=256
\]
个磁盘块号。

所以一次间接地址覆盖的逻辑块号为
\[
8\sim(8+256-1)=8\sim263.
\]

逻辑块 269 已经超过 263，因此需要使用\textbf{二次间接地址}。

在已经获得 f.dat 的 I 节点地址信息之后，为读到逻辑块 269，需要：
\begin{enumerate}
  \item 读二次间接的第一级间接块，取得第二级间接块号；
  \item 读第二级间接块，取得实际数据块号；
  \item 读实际数据块。
\end{enumerate}

因此共需
\[
\boxed{3}\text{ 个磁盘块}.
\]
\end{explanation}

\section*{五、（12分）调度方案分析。}

\begin{question}
某应用场景分 2 类任务：交互式任务和后台任务，要求交互式任务绝对优先于后台任务执行，交互式任务之间、后台任务之间没有优先级的区别。小明为该应用场景的专用操作系统设计了一个调度方案：设置 6 个就绪队列 Q[6]，i（0$\leq$i<6）越小，Q[i] 的级别越高，调度程序总是选择执行级别高的队首任务，只有 Q[0]-Q[i-1] 为空，才会调度 Q[i] 的任务；就绪的后台任务进入 Q[5] 的队尾，就绪的交互式任务进入 Q[0] 的队尾，每个队列采用时间片轮转调度，执行的时间片到时被剥夺，所有队列的时间片大小相同；当高级队列有任务就绪时，正在执行的低级队列的任务被剥夺。交互式任务被剥夺时，降到下一级队列的队尾（Q[4] 的任务不降级）；后台任务被剥夺时，排在 Q[5] 的队尾。

1. 对于交互式任务，小明的设计方案存在饥饿的问题，请说明为什么？给出一种解决该问题的方法，并加以解释。

2. 对于交互式任务，在排除饥饿的情况下，小明的设计方案存在不够公平的问题，请说明为什么？给出一种解决该问题的方法，并加以解释。

3. 对于后台任务，给出一种与小明的设计相比调度开销更小的方法，解释理由，并说明你的方法可能带来的其它方面的不足。
\end{question}

\begin{answer}
1. 低级交互式队列可能因高级队列持续有任务而长期得不到调度，产生饥饿。可采用老化或周期性把任务提升到高优先级队列。

2. 同属交互式任务却被放入不同优先级队列，新任务长期优先于已降级任务，不够公平。可把所有交互式任务统一放入一个高优先级 RR 队列。

3. 后台任务可采用 FCFS，仅在交互式任务到来时被抢占，从而减少时间片到期引起的上下文切换。缺点是长后台任务可能使其它后台任务等待过久，公平性和响应时间较差。
\end{answer}

\begin{explanation}
\textbf{第 1 问：饥饿。}

交互式任务刚就绪时进入 Q[0]，以后每次因为时间片到或受到更高级队列任务影响而被剥夺，都可能逐级下降到 Q[1]、Q[2]、Q[3]、Q[4]。调度器永远先选择高优先级队列。

如果不断有新的交互式任务进入 Q[0]，或者高优先级队列长期保持非空，那么已经下降到 Q[3]、Q[4] 的旧任务就可能长期得不到 CPU，即发生饥饿。

一种典型解决方法是\textbf{老化（aging）/周期性优先级提升}：任务等待时间越长，逐步提高它的优先级；或者每隔一定时间把所有交互式任务统一提升回 Q[0]。这样任何等待足够久的交互式任务最终都会进入较高优先级，从而有机会执行。

\medskip
\textbf{第 2 问：公平性。}

题目前提说明“交互式任务之间没有优先级的区别”，但该方案却因为任务到达时间和被剥夺次数不同，把交互式任务放在 Q[0]--Q[4] 的不同等级中。

例如，一个刚到达的交互式任务直接进入 Q[0]，而一个已经运行过、被剥夺过多次的交互式任务可能处于 Q[4]。即使已经通过老化等办法保证后者最终不会饿死，新任务仍然会在相当长的时间里优先于旧任务执行。这样，同属交互式任务的进程获得 CPU 的机会并不平等。

一种更直接的改进是：\textbf{所有交互式任务只使用一个高优先级就绪队列，并在该队列内采用时间片轮转。}所有交互式任务具有相同优先级，按照 RR 轮流执行；只要交互式队列非空，就不调度后台任务，因此仍然满足“交互式任务绝对优先于后台任务”的要求。

也可以保留多级队列，但规定只有真正用完整个时间片的任务才降级，而因更高级任务到达而提前被抢占的任务保持原级，并配合周期性优先级提升。不过从本题“交互式任务之间没有优先级区别”的要求看，统一 RR 队列更直接。

\medskip
\textbf{第 3 问：降低后台任务调度开销。}

后台任务之间没有优先级差别，而且只在没有交互式任务时才执行。可以把后台任务放在一个\textbf{FCFS 队列}中：后台任务一旦获得 CPU，就一直运行到结束或阻塞；只有出现交互式任务时才立即抢占后台任务。

与固定小时间片的 RR 相比，后台任务之间不再因为时间片耗尽而频繁发生上下文切换，也减少了定时器中断和重新排队的开销，因此调度开销更小。

不足是：如果某个后台任务 CPU 执行时间很长，它可能长时间占用后台任务可获得的 CPU，使排在后面的后台任务等待较久，响应时间和公平性都会变差，容易出现“护航效应（convoy effect）”。
\end{explanation}

\section*{六、（18分）同步算法设计。}

\begin{question}
某银行的业务分为现金和转账，这 2 项业务分别由 2 个和 1 个服务员负责办理。业务大厅最多能容纳 20 位等候办理业务的客户（不包括正在办理业务的客户）。客户到来时，若业务大厅已满，则离开；否则进入大厅根据办理的业务取号，然后等候服务员叫号，办理完业务后离开。试用信号量的 P、V 操作设计该问题的同步算法，给出所用共享变量（如果需要）和信号量及其初始值，并说明各自的含义。
\end{question}

\begin{answer}
一种可行设计：
\[
mutex=1,\qquad cash=2,\qquad transfer=1,
\]
共享变量 $waiting=0$。

客户到达后先在 $mutex$ 保护下检查 $waiting$：若等于 20 则立即离开；否则 $waiting++$ 并取号。随后现金客户执行 $P(cash)$，转账客户执行 $P(transfer)$。成功获得服务员后立即令 $waiting--$，办理完业务后分别执行 $V(cash)$ 或 $V(transfer)$。

若要求严格按号叫号，可让相应等待队列采用 FIFO。
\end{answer}

\begin{explanation}
这道题有三个关键约束：
\begin{enumerate}
  \item 大厅最多只能有 20 个\textbf{正在等待}的客户；已经开始办理业务的客户不计入这 20 人。
  \item 现金业务同时最多有 2 人正在办理，因为有 2 个现金服务员。
  \item 转账业务同时最多有 1 人正在办理，因为只有 1 个转账服务员。
\end{enumerate}

可以使用一个共享变量记录大厅内当前等待人数，再用互斥信号量保护它。两个“服务员资源信号量”分别表示当前可用的现金服务员和转账服务员。

\textbf{共享变量：}
\begin{lstlisting}
int waiting = 0;        // 当前在大厅等待、尚未开始办理业务的人数
int cashNo = 0;         // 现金业务取号计数（如需要显示号码）
int transferNo = 0;     // 转账业务取号计数（如需要显示号码）
\end{lstlisting}

\textbf{信号量及初值：}
\begin{align*}
mutex &= 1 &&\text{保护 waiting 以及取号计数器等共享变量；}\\
cash &= 2 &&\text{表示 2 个现金服务员；}\\
transfer &= 1 &&\text{表示 1 个转账服务员。}
\end{align*}

客户进程可以写成下面的形式。

\begin{lstlisting}
Customer(type)
{
    P(mutex);
    if (waiting == 20)
    {
        V(mutex);
        leave();                // 大厅已满，直接离开
        return;
    }

    waiting++;                  // 进入大厅，成为等待客户

    if (type == CASH)
        myNo = ++cashNo;        // 现金业务取号
    else
        myNo = ++transferNo;    // 转账业务取号

    V(mutex);

    if (type == CASH)
        P(cash);                // 等待任意一个现金服务员叫号
    else
        P(transfer);            // 等待转账服务员叫号

    // P 成功，说明已经轮到该客户，开始办理业务。
    // 此时他不再属于“大厅中等待的人”。
    P(mutex);
    waiting--;
    V(mutex);

    do_service(type);           // 办理业务

    if (type == CASH)
        V(cash);                // 现金服务员空闲，可叫下一位
    else
        V(transfer);            // 转账服务员空闲，可叫下一位

    leave();
}
\end{lstlisting}

如果要求严格按照号码先后顺序叫号，可以规定现金业务和转账业务各自的信号量等待队列采用 FIFO，或者分别维护一个按取号顺序排列的 FIFO 等待队列。这样，被唤醒的客户顺序就与取号顺序一致。

这里最容易写错的是 $waiting$ 的减少时机：客户一旦通过 $P(cash)$ 或 $P(transfer)$ 获得服务员，说明已经开始办理业务，因此应当马上执行 $waiting--$。不能等办理完离开时才减，否则就把“正在办理业务的客户”错误地算进大厅 20 人的容量中。
\end{explanation}
\end{document}
